% 第二阶段补充材料：推导链路与数值实现。AI 起草，逐阶段经用户过目后写入。
% 文件名 ASCII（本机 biber 不认中文文件名）；中文主题见 notes/README.md 索引。
\documentclass[UTF8,11pt,a4paper]{ctexart}
\usepackage[margin=2.6cm]{geometry}
\usepackage{amsmath,amssymb}
\usepackage{booktabs}
\usepackage{longtable}
\usepackage[colorlinks=true,linkcolor=blue,citecolor=blue,urlcolor=blue]{hyperref}
\usepackage[backend=biber,style=numeric,sorting=none,doi=true,url=false,eprint=true]{biblatex}
\renewcommand*{\bibfont}{\small}
\setlength{\bibitemsep}{0pt}
\addbibresource{../../literature/references.bib}

\title{第二阶段补充材料：推导链路与数值实现\\[0.4em]
\large 从模型建立到热三临界候选与相图拓扑的完整过程}
\author{Quantum Optics Thesis 工作笔记}
\date{2026 年 11 月 27 日}

\begin{document}
\maketitle

\begin{abstract}
本材料是[第二阶段研究汇总]的推导补充：正文面向汇报、给出结论与证据分级，本材料
面向"每个结论是如何一步步得到的"。按阶段展开系统选择、模型建立、热力学自洽、
相边界判据、零场 Landau 系数、零温校准、数值方法、热三临界候选与相图拓扑，每个阶段
区分理论推导、数值实现与所得物理内容。符号统一见阶段 0 的约定表。
\end{abstract}

\section{阶段 0：系统选择与模型建立}

\subsection{研究对象与轴向选择}

我们研究的是一维\textbf{量子 Dicke--Ising 链}：$N$ 个自旋 $1/2$ 排成周期链，自旋之间
最近邻 Ising 交换、外加静态纵场，并集体耦合到单个腔模。用三元组
$(\text{交换轴}/\text{静场轴}/\text{腔耦合轴})$ 标记轴向，本项目取
\textbf{$zz/z/x$}：交换与静场都沿 $z$（同轴），腔耦合沿 $x$（横向）。

选这个轴向由三点决定：

\begin{enumerate}
  \item \textbf{三轴夹角是不变量。} 全局自旋旋转同时转动三个轴，不能改变任意两轴的
        夹角。$zz/z/x$ 中交换轴与静场轴\textbf{重合}（夹角 $0^\circ$），而三轴两两
        正交的 $yy/z/x$（Gammelmark--Mølmer 2011）与 $zz/y/x$（Rao 2025）夹角是
        $90^\circ$，二者严格等价，且已被有限温一/二级边界与三临界研究覆盖。我们选的
        是未被覆盖的另一等价类。
  \item \textbf{零场极限的物理不同。} $h=0$ 时，本模型的物质链是
        $J\sigma^z\sigma^z+(\Delta/2)\sigma^z$，两项对易、纯对角，是\textbf{经典}
        Ising 链；而 $yy/z/x$ 是 $J\sigma^y\sigma^y+(\Delta/2)\sigma^z$，不对易，是
        \textbf{量子}横场 Ising 链。这决定可积性分界与求解方法。
  \item \textbf{平台可实现性。} 交换轴由相互作用的各向异性决定、最难调，静场轴由外加
        DC 场决定；让二者同轴是磁偶极原子、Rydberg 修饰、超导比特、磁振子晶体等平台的
        默认配置。（详细论证见轴向等价类与平台对应说明，标注为待核推断。）
\end{enumerate}

\subsection{哈密顿量}

取自然单位 $\hbar=k_B=1$、自旋 $s_i^\alpha=\sigma_i^\alpha/2$（$\sigma^\alpha$ 为泡利
矩阵）。完整哈密顿量为
\begin{equation}
  H_N=\omega a^\dagger a+\Delta\sum_i s_i^z+4J\sum_i s_i^z s_{i+1}^z
  +\frac{2g}{\sqrt N}(a+a^\dagger)\sum_i s_i^x .
  \label{eq:HN}
\end{equation}
四项依次是腔模能量（$\omega>0$）、静态纵场（自旋能级劈裂/失谐）、最近邻 Ising 交换
（$J<0$ 铁磁、$J>0$ 反铁磁）、腔--自旋耦合（沿 $x$，$1/\sqrt N$ 保证热力学极限良好
定义）。本项目主要取铁磁 $J=-0.2$、$\omega=1$；纵场 $\Delta>0$ 可经全局 $z$ 翻转选为正。

\subsection{符号约定}

下表是后续所有推导的统一依据，改符号时须同步更新本表与相关公式。

{\footnotesize
\begin{longtable}{@{}p{2.1cm}p{6.0cm}p{6.2cm}@{}}
\toprule
符号 & 含义 & 约定 / 取值 \\
\midrule
\endhead
\multicolumn{3}{l}{\textbf{基础}} \\
$N$ & 自旋（位点）数 & 周期链，$N\to\infty$ 为热力学极限 \\
$\sigma_i^\alpha,\ s_i^\alpha$ & 泡利矩阵 / 自旋算符 & $s_i^\alpha=\sigma_i^\alpha/2$，$\alpha=x,y,z$ \\
$\omega$ & 腔模频率 & $\omega>0$，取 $\omega=1$ \\
$a,\ a^\dagger$ & 腔模湮灭 / 产生算符 & 玻色子 \\
$\Delta$ & 纵场（失谐） & $\Delta>0$ 可经全局 $z$ 翻转选定 \\
$J$ & 最近邻 Ising 交换 & $J<0$ 铁磁（取 $-0.2$），$J>0$ 反铁磁 \\
$g$ & 腔--自旋耦合强度 & 单模集体耦合 \\
$T,\ \beta$ & 温度 / 逆温 & $\beta=1/T$ \\
\midrule
\multicolumn{3}{l}{\textbf{自洽量}} \\
$\varphi$ & 腔位移序参量 & $\varphi=\mathrm{Re}\langle a\rangle/\sqrt N$，实数 \\
$h$ & 等效横场 & $h=4g\varphi\ge0$ \\
$f_{\rm spin}(T,h)$ & 冻结腔后的无限链单位点自旋自由能 & $f_{\rm spin}=-\lim_{N\to\infty}(T/N)\log\operatorname{Tr}e^{-H_{\rm spin}(h)/T}$ \\
$\Phi_T(\varphi)$ & 总自由能泛函 & $\Phi_T(\varphi)=\omega\varphi^2+f_{\rm spin}(T,4g\varphi)$ \\
$A_T(h)$ & 正常态与有限场自由能差 & $A_T(h)=f_{\rm spin}(T,0)-f_{\rm spin}(T,h)$ \\
\midrule
\multicolumn{3}{l}{\textbf{判据}} \\
$g_{\rm equal}(h)$ & 正常态与有限位移态等能耦合 & $g_{\rm equal}^2=\omega h^2/[16A_T(h)]$ \\
$g_{\rm onset}(T)$ & 等能耦合下确界 & $g_{\rm onset}=\inf_{h>0}g_{\rm equal}(h)$ \\
$g_{\rm local}(T)$ & 正常态局部失稳阈值 & 由零场曲率决定（阶段 2） \\
$Q$ & 三自旋概率加权因子 & 见阶段 2 \\
$\chi_x$ & 横向磁化率 & $\chi_x=Q/4$ \\
$g_c(\Delta)$ & 零温一级线 & 全局等能最小值（阶段 4） \\
\midrule
\multicolumn{3}{l}{\textbf{Landau 系数}} \\
$a_2,a_4,a_6$ & $f_{\rm spin}(T,h)$ 在 $h=0$ 的展开系数 & $f=f_0+a_2h^2+a_4h^4+a_6h^6+\cdots$ \\
$Z(h),\ Z_m$ & 配分函数及其展开 & $Z(h)=Z_0+h^2Z_2+h^4Z_4+h^6Z_6+\cdots$，$z_m=Z_m/Z_0$ \\
$dE_i$ & 翻转单自旋的能量差 & 见阶段 3 \\
$dE_{ij}$ & 翻转相邻两自旋的能量差 & $dE_{ij}=dE_i+dE_j+4Js_is_j$ \\
$B_1,\ I_4,\ F_m$ & 虚时间积分核 / 合流差商 & 见阶段 3 \\
\midrule
\multicolumn{3}{l}{\textbf{零温基线}} \\
$K,\ \Gamma$ & 横场 Ising 的交换 / 横场 & $K=-J>0$，$\Gamma=h/2$ \\
$A,\ B$ & 零温展开辅助量 & $A=\Delta/2-2J$，$B=\Delta/2-J$ \\
$e_{\rm spin}(0,h)$ & 零温基态能密度 & 见阶段 4 \\
\midrule
\multicolumn{3}{l}{\textbf{证据量}} \\
$h^\ast$ & 一级前驱极小位置 & $h^{\ast\,2}=-a_4/(2a_6)$ \\
$\mathrm{depth}$ & 一级前驱极小深度 & $\mathrm{depth}=g_{\rm local}-g_{\rm equal}(h^\ast)$ \\
$T^\ast,\ g^\ast$ & 热三临界候选 & $a_4(T^\ast)=0$，$g^\ast=g_{\rm local}(T^\ast)$ \\
$T_{\rm upper},\ T_{\rm lower}$ & 上支 / 下支 & $a_4(T)=0$ 的高温根 / 低温根 \\
$\Delta_{\rm end},T_{\rm end},g_{\rm end}$ & 切触端点 & 上下支汇合点 \\
\bottomrule
\end{longtable}
}

\end{document}
